\section{Conclusion}
\label{sec:conclusion}

We measure escape-time distributions of SGD and Adam from exact minima and saddles of small tanh networks, and we test four numerical claims against criteria fixed in advance.
Escape from these basins is a discrete hop with a non-constant hazard, shaped mainly by a delay within a minimum and by large differences between minima, not by a Weibull law with $k$ in 0.7--1.3 tied to geometry.
$\etaB$ is the right control variable on the main landscape, but $\tau$ scales about as $B^{2}$, not $\sqrt{B}$.
Adam's advantage is real at matched in-basin loss (about $2\times$, from adaptive scaling and not momentum), is not shown to be specific to narrow basins, and disappears at the default learning rate.
Saddle exit follows $\etop$ at low noise, the gap does not predict how often, and at high noise exit is reproducible along other directions.

The main takeaway is that no single shape parameter summarizes escape: delay, minimum identity and noise level each change the answer.
Future work should extend run lengths at large $B$ to settle power law versus Arrhenius, start from conditioned quasi-stationary states to measure the delay directly, test the non-$\etop$ saddle exits for noise-induced drift, and repeat these measurements on wider or deeper networks with Hessian-vector quenching.
One open question remains: why does the hazard fall late within a single minimum, through competing channels or slow drift to harder regions?
